%% --------------------------------------------------------
\section{Метод DIIS та прискорення збіжності SCF}\label{sec:DIIS}
%% --------------------------------------------------------

У стандартному SCF-алгоритмі ми на кожній ітерації будуємо нову матрицю Фока $\vect{F}^{(k)}$ з поточної матриці густини $\vect{D}^{(k-1)}$, розв'язуємо задачу на власні значення та отримуємо нову матрицю густини $\vect{D}^{(k)}$. Цей процес повторюється до збіжності.

Проблема у тому, що такий простий ітераційний процес може збігатися дуже повільно, особливо для складних систем. Іноді енергія навіть осцилює між кількома значеннями і не досягає стабільного результату.

Метод \textbf{DIIS} (Direct Inversion in the Iterative Subspace --- пряма інверсія в ітераційному підпросторі), запропонований Пітером Пулаєм у 1980 році, вирішує цю проблему наступним чином: замість того, щоб використовувати лише матрицю Фока з поточної ітерації, DIIS \emph{запам'ятовує} кілька попередніх матриць Фока
\begin{equation*}
	\vect{F}^{(n-k)}, \vect{F}^{(n-k+1)}, \ldots, \vect{F}^{(n-1)}, \vect{F}^{(n)}
\end{equation*}
та будує оптимальну лінійну комбінацію:

\begin{equation}
	\vect{F}_{\text{DIIS}}^{(n)} = \sum_{i=1}^{m} c_i \vect{F}^{(n-m+i)}
\end{equation}
%
де коефіцієнти $c_i$ підбираються так, щоб мінімізувати похибку самоузгодженості.

\begin{commentbox}[Що таке похибка самоузгодженості?]
	Якщо рішення ідеально самоузгоджене, то виконується комутаційне співвідношення:
	\begin{equation}
		[\vect{F}, \vect{D}\vect{S}] = \vect{F}\vect{D}\vect{S} - \vect{S}\vect{D}\vect{F} = \vect{0}
	\end{equation}

	Похибка (residual) на $i$-тій ітерації --- це відхилення від цієї умови:
	\begin{equation}
		\vect{R}^{(i)} = \vect{F}^{(i)} \vect{D}^{(i)} \vect{S} - \vect{S} \vect{D}^{(i)} \vect{F}^{(i)}
	\end{equation}

	DIIS шукає такі коефіцієнти $c_i$, щоб похибка комбінованої матриці Фока була мінімальною.
\end{commentbox}

SCF-процес умовно можна уявити як <<рух по енергетичному ландшафту>> до мінімуму. Без DIIS процедура робить кроки, які можуть <<перестрибувати>> через мінімум туди-сюди. DIIS аналізує кілька попередніх кроків і робить <<розумну>> екстраполяцію, яка швидше веде до мінімуму.

Типова схема роботи:
\begin{enumerate}
	\item \textbf{Ітерації 1--2:} звичайний SCF без DIIS (потрібно накопичити історію)
	\item \textbf{Ітерація 3 і далі:} DIIS активується автоматично, зберігає 6--8 останніх матриць Фока та похибок
	\item На кожній ітерації будується оптимальна комбінація, яка прискорює збіжність
	\item Процес продовжується до досягнення заданої точності (наприклад, $\Delta E < 10^{-6}$ Hartree)
\end{enumerate}

Як керувати \texttt{DIIS}, та що робити, коли він не працює, наведено в додатку~\ref{sec:convergence}.
